\section{Risk Assessment}
% \begin{itemize}
%     \item What are the potential risks related to your project?
%     \item What issues may you expect during development?
%     \item How would you circumvent this?
%     \item How likely is this project to be successful?
%     \item What are societal or user risks when your model is deployed publicly?
%     \item \textbf{Required:} At least one threat to your project development and an estimate of success.
%     \item \textbf{Required:} A consideration of possible risks of having your model deployed.
% \end{itemize}

One major risk is the use of live streaming data, since this is sensitive information. For that reason, we would prefer to keep the system on-device rather than stream sensor data to a server.

We may also face a trade-off between model size and accuracy, and we may need to learn deployment techniques that are specific to this use case. To prepare for that, we want to look early at projects such as the web \texttt{fake beer} drinking app to understand how live streaming sensor data is handled, and we may also need to explore mobile development.

The project is likely to be successful in terms of model accuracy, since the task has already been shown to be possible with high accuracy. What is less certain is how feasible the approach will be once compression and deployment constraints are added.

The main societal risk is misuse: the system could be used to detect, follow, or identify someone through their embedding, which would effectively leak biometric information. There are also user risks because this system would be an added security layer, and users might blame the model if it fails in a sensitive situation. In a real deployment, the model would need to achieve a very high level of reliability before it could be trusted.